\documentclass[11pt,a4paper]{article}
\usepackage[utf8]{inputenc}
\usepackage[T1]{fontenc}
\usepackage{amsmath,amsfonts,amssymb}
\usepackage{graphicx}
\usepackage{hyperref}
\usepackage{listings}
\usepackage[dvipsnames]{xcolor}
\usepackage{booktabs}
\usepackage{float}
\usepackage{geometry}
\geometry{margin=1in}

\definecolor{zenblue}{RGB}{41,121,255}
\definecolor{zengreen}{RGB}{52,199,89}
\definecolor{codegray}{RGB}{245,245,245}

\hypersetup{colorlinks=true,linkcolor=zenblue,urlcolor=zenblue,citecolor=zenblue}

\lstset{
    backgroundcolor=\color{codegray},
    basicstyle=\ttfamily\small,
    breaklines=true,
    captionpos=b,
    frame=single,
    numbers=left,
    numberstyle=\tiny\color{gray}
}

\title{
    \vspace{-2cm}
    \Large \textbf{Zen AI Model Family} \\
    \vspace{0.5cm}
    \Huge \textbf{Zen-World} \\
    \vspace{0.3cm}
    \large Packaging Wan2.1 Text-to-Video Generation for the Zen Stack \\
    \vspace{0.5cm}
    \normalsize Technical Report
}

\author{
    Hanzo AI Research Team\thanks{research@hanzo.ai} \and
    Zoo Labs Foundation\thanks{foundation@zoo.ngo}
}

\date{June 2026}

\begin{document}

\maketitle

\begin{abstract}
\textbf{Zen-World} is a packaging and integration of \textbf{Wan2.1-T2V-14B}, the open-source
text-to-video diffusion model developed and released by Alibaba's Team Wan, into the Zen model
stack. Zen-World is \emph{not} a from-scratch model: it redistributes the upstream
\texttt{Wan-AI/Wan2.1-T2V-14B} weights and architecture under the permissive Apache~2.0 license,
adding Zen-stack loaders, configuration, and a unified inference interface. The underlying model is
a $\sim$14-billion-parameter Diffusion Transformer (DiT) trained with a flow-matching objective,
paired with the Wan2.1 3D causal variational autoencoder (Wan-VAE) for spatio-temporal latent
compression and the multilingual umT5 text encoder for prompt conditioning. It generates 480P and
720P video from English or Chinese text prompts. This report describes the real architecture as
released upstream, attributes it to its authors, and documents how it is integrated into Zen. All
architectural and capability claims here are those of the upstream Wan2.1 release; we do not report
any independent benchmark numbers of our own.
\end{abstract}

\section*{Attribution and License}
\addcontentsline{toc}{section}{Attribution and License}

\textbf{Upstream model.} Zen-World redistributes \texttt{Wan-AI/Wan2.1-T2V-14B}, part of the
\emph{Wan} series of open large-scale video generative models from Alibaba's Team Wan
\cite{wan2025report}.

\textbf{License.} The Wan2.1 models and source code are released under the
\textbf{Apache License 2.0}. Zen-World preserves this license and the upstream copyright and
attribution notices. No relicensing is claimed; redistribution is permitted because Apache~2.0 is a
permissive license.

\textbf{What Zen adds.} Zen-World contributes \emph{packaging only}: Zen-stack model loaders, a
consistent configuration and weights layout, and a unified text-to-video inference API consistent
with the rest of the Zen family. It does not modify the model architecture or retrain the weights.

\tableofcontents
\newpage

\section{Introduction}

Text-to-video generation synthesizes a short video clip from a natural-language description. Modern
systems follow the diffusion-transformer paradigm: a transformer denoiser operates in a compressed
video latent space produced by a video autoencoder, conditioned on text embeddings from a language
encoder, and is trained to reverse a noising (or flow-matching) process \cite{wan2025report}.

Zen-World packages one such system---\textbf{Wan2.1-T2V-14B}---into the Zen stack. Wan2.1 was
released by Alibaba's Team Wan as an open, Apache-2.0-licensed family of video generative models
\cite{wan2025report}. The 14B text-to-video variant is the model redistributed here. Zen-World
provides Zen-native loading and a unified inference interface; the generative capability, weights,
and architecture are entirely those of the upstream release.

This report (1) attributes the model to its authors and states its license, (2) describes the real
Wan2.1 architecture and its components, (3) clarifies the scope of Zen's integration, and
(4) reports only capabilities and figures stated by the upstream release, with attribution.

\subsection{Model Overview}

\begin{table}[H]
\centering
\begin{tabular}{ll}
\toprule
\textbf{Property} & \textbf{Value} \\
\midrule
Upstream model & \texttt{Wan-AI/Wan2.1-T2V-14B} (Alibaba Team Wan) \\
Task & Text-to-video generation \\
License & Apache 2.0 \\
Architecture & Diffusion Transformer (DiT), flow matching \\
Denoiser parameters & $\sim$14B \\
Latent autoencoder & Wan2.1 3D causal VAE (Wan-VAE) \\
Text encoder & umT5 (multilingual; English / Chinese prompts) \\
Output resolutions & 480P ($832\times480$), 720P ($1280\times720$) \\
Zen contribution & Packaging / integration into the Zen stack \\
\bottomrule
\end{tabular}
\caption{Zen-World summary. Architecture and capability values are those of upstream Wan2.1.}
\end{table}

\section{Architecture}

Wan2.1 follows the mainstream diffusion-transformer design and is composed of three primary
components \cite{wan2025report}: a 3D causal VAE (\emph{Wan-VAE}), a multilingual umT5 text encoder,
and a flow-matching Diffusion Transformer (DiT) denoiser. Figure-level overview:

\begin{enumerate}
    \item \textbf{Wan-VAE} encodes video into a compact spatio-temporal latent and decodes
    generated latents back to pixels.
    \item \textbf{umT5} encodes the text prompt into embeddings.
    \item \textbf{DiT} denoises the video latent, conditioned on the text embeddings via
    cross-attention, under a flow-matching objective.
\end{enumerate}

\subsection{Wan-VAE: 3D Causal Video Autoencoder}

Generating video directly in pixel space is prohibitively expensive, so Wan2.1 operates in a
compressed latent space produced by \textbf{Wan-VAE}, a 3D \emph{causal} variational autoencoder
designed for video \cite{wan2025report}. The VAE performs spatio-temporal compression with a
temporal downsampling factor of $t=4$ and spatial downsampling factors of $h=w=8$, while preserving
temporal causality so that motion remains smooth across frames. The latent has 16 channels (the
DiT's input and output channel dimension; see Table~\ref{tab:dit}). The causal design lets the VAE
encode and decode long videos without breaking temporal consistency.

\subsection{umT5 Text Encoder}

Prompts are encoded by \textbf{umT5}, a multilingual T5-family text encoder. The resulting text
embeddings condition the DiT through cross-attention. The multilingual encoder is one reason Wan2.1
accepts both English and Chinese prompts, and---per the upstream release---Wan2.1 is described as
the first video model able to generate both Chinese and English text within the rendered video
\cite{wan2025report}. The text sequence length used by the 14B configuration is 512 tokens
(Table~\ref{tab:dit}).

\subsection{Diffusion Transformer (DiT)}

The denoiser is a Diffusion Transformer. It processes the noised video latent as a sequence of
patch tokens through a stack of transformer blocks. Each block (the \emph{WanAttentionBlock})
combines, per the upstream design, vision self-attention over the video tokens, text-to-vision
cross-attention that injects the umT5 prompt embeddings, and a feed-forward layer, with timestep
conditioning applied to modulate the block \cite{wan2025report}.

The exact configuration of the redistributed 14B text-to-video model, taken directly from the
upstream model configuration, is given in Table~\ref{tab:dit}.

\begin{table}[H]
\centering
\begin{tabular}{ll}
\toprule
\textbf{Hyperparameter} & \textbf{Value (T2V-14B)} \\
\midrule
Model type & \texttt{t2v} (text-to-video) \\
Model / embedding dimension (\texttt{dim}) & 5120 \\
Transformer layers (\texttt{num\_layers}) & 40 \\
Attention heads (\texttt{num\_heads}) & 40 \\
Feed-forward dimension (\texttt{ffn\_dim}) & 13824 \\
Latent input / output channels (\texttt{in\_dim} / \texttt{out\_dim}) & 16 \\
Frequency embedding dimension (\texttt{freq\_dim}) & 256 \\
Text sequence length (\texttt{text\_len}) & 512 \\
LayerNorm epsilon (\texttt{eps}) & $1\times10^{-6}$ \\
\bottomrule
\end{tabular}
\caption{Wan2.1-T2V-14B DiT configuration, from the upstream \texttt{config.json}.}
\label{tab:dit}
\end{table}

\subsection{Flow-Matching Objective}

Rather than a traditional discrete-noise diffusion schedule, Wan2.1 is trained with
\textbf{flow matching} \cite{lipman2023flow, wan2025report}. Flow matching defines intermediate
states by linear interpolation between a sample and noise and trains the network to predict the
corresponding velocity (the time derivative of the interpolation path). Concretely, for a clean
latent $x_1$ and noise $x_0 \sim \mathcal{N}(0, I)$, an intermediate latent is
\begin{equation}
    x_t = (1 - t)\, x_0 + t\, x_1, \qquad t \in [0, 1],
\end{equation}
and the model $v_\theta$ is trained to regress the target velocity $x_1 - x_0$:
\begin{equation}
    \mathcal{L}_{\text{FM}} = \mathbb{E}_{t,\, x_0,\, x_1,\, c}
    \left\lVert v_\theta(x_t, t, c) - (x_1 - x_0) \right\rVert_2^2,
\end{equation}
where $c$ denotes the umT5 text conditioning. At inference, video latents are generated by
integrating the learned velocity field from noise, then decoded to pixels by Wan-VAE.

\section{Capabilities (as reported upstream)}

The capabilities below are those stated by the upstream Wan2.1 release \cite{wan2025report}; they
are reproduced here with attribution. Zen-World does not add or independently re-measure these
figures.

\begin{itemize}
    \item \textbf{Output.} Generates short video clips at 480P ($832\times480$) and 720P
    ($1280\times720$) from text prompts.
    \item \textbf{Multilingual prompts and in-video text.} Accepts English and Chinese prompts;
    described upstream as the first video model able to render both Chinese and English text in the
    generated video.
    \item \textbf{Model family.} The Wan2.1 release also includes a smaller T2V-1.3B model and
    image-to-video (I2V), first-last-frame (FLF2V), and video-editing (VACE) variants; Zen-World
    redistributes the T2V-14B variant.
    \item \textbf{Efficiency (sibling model).} The upstream release reports that the smaller
    T2V-1.3B model can generate a 5-second 480P clip using about 8.19~GB of VRAM, enabling
    consumer-GPU use; the 14B model redistributed here is correspondingly heavier.
    \item \textbf{Evaluation.} The upstream report states the 14B model was evaluated on an internal
    suite of 1{,}035 prompts across 14 dimensions and that it is competitive with leading
    open-source and commercial systems. We reproduce this claim as reported and do not restate
    specific scores here.
\end{itemize}

\section{Integration into the Zen Stack}

Zen-World's contribution is integration, not modeling. Specifically, Zen-World provides:
\begin{itemize}
    \item Zen-stack loading of the upstream Wan2.1-T2V-14B weights (DiT, Wan-VAE, umT5).
    \item A consistent on-disk configuration and weights layout aligned with the Zen family.
    \item A unified text-to-video inference interface consistent with other Zen models.
\end{itemize}
No architectural changes are made and the weights are not retrained. The example below illustrates
the intended Zen-stack interface for text-to-video generation.

\begin{lstlisting}[language=Python, caption=Zen-World text-to-video (packaging of Wan2.1-T2V-14B)]
from zen import ZenWorld

# Loads the redistributed Wan2.1-T2V-14B weights (DiT + Wan-VAE + umT5).
model = ZenWorld.from_pretrained("zenlm/zen-world")

# Text-to-video generation.
video = model.generate(
    prompt="A red panda walking through a snowy bamboo forest, cinematic lighting",
    resolution="720p",      # 1280x720; also "480p" (832x480)
    num_frames=81,
    fps=16,
)

video.save("output.mp4")
\end{lstlisting}

\section{Related Work}

Zen-World is a redistribution of Wan2.1 \cite{wan2025report}, which belongs to the broader line of
diffusion-transformer video generators. The diffusion-transformer (DiT) backbone derives from the
scalable transformer denoiser of Peebles and Xie \cite{peebles2023dit}. Flow matching, the training
objective Wan2.1 uses in place of a discrete noise schedule, was introduced by Lipman et al.
\cite{lipman2023flow}. The Wan technical report \cite{wan2025report} details the Wan-VAE 3D causal
autoencoder, the umT5 conditioning, and the model family.

\section{Conclusion}

Zen-World is the Zen stack's packaging of \textbf{Wan2.1-T2V-14B}, the Apache-2.0 text-to-video
diffusion model released by Alibaba's Team Wan. The real system is a $\sim$14B Diffusion Transformer
trained with flow matching, operating in the latent space of the Wan2.1 3D causal VAE and
conditioned by the multilingual umT5 text encoder, producing 480P and 720P video from English or
Chinese prompts. Zen's contribution is integration: loaders, configuration, and a unified inference
interface. All architecture and capability statements in this report are attributed to the upstream
Wan2.1 release, and no independent benchmark numbers are claimed.

\begin{thebibliography}{10}
\bibitem{wan2025report} Team Wan (Alibaba), ``Wan: Open and Advanced Large-Scale Video Generative Models,'' arXiv:2503.20314, 2025. Models and code: \url{https://github.com/Wan-Video/Wan2.1}; weights: \url{https://huggingface.co/Wan-AI/Wan2.1-T2V-14B} (Apache~2.0).
\bibitem{lipman2023flow} Y. Lipman, R. T. Q. Chen, H. Ben-Hamu, M. Nickel, and M. Le, ``Flow Matching for Generative Modeling,'' ICLR, 2023. arXiv:2210.02747.
\bibitem{peebles2023dit} W. Peebles and S. Xie, ``Scalable Diffusion Models with Transformers,'' ICCV, 2023. arXiv:2212.09748.
\end{thebibliography}

\end{document}
